\documentclass[11pt,a4paper]{article}
\usepackage[utf8]{inputenc}
\usepackage{amsmath,amssymb,amsthm}
\usepackage[margin=1in]{geometry}

\newtheorem{theorem}{Theorem}
\newtheorem{lemma}[theorem]{Lemma}
\newtheorem{proposition}[theorem]{Proposition}
\newtheorem{definition}[theorem]{Definition}

\title{Problem 122: Efficient Exponentiation}
\author{Project Euler Solution}
\date{}

\begin{document}
\maketitle

\section{Problem Statement}

Find $\sum_{n=1}^{200} m(n)$ where $m(n)$ is the length of a shortest addition chain for $n$.

\section{Mathematical Development}

\begin{definition}[Addition chain]
An addition chain for $n$ is a sequence $1 = a_0 < a_1 < \cdots < a_r = n$ where each $a_k = a_i + a_j$ for some $0 \le i \le j < k$. The length is $r$, and $m(n)$ denotes the minimum such $r$.
\end{definition}

\begin{theorem}[Lower bound]
For all $n \ge 1$, $m(n) \ge \lceil \log_2 n \rceil$.
\end{theorem}

\begin{proof}
By induction, $a_k \le 2^k$ (each step at most doubles the maximum). Hence $n = a_r \le 2^r$, giving $r \ge \lceil \log_2 n \rceil$.
\end{proof}

\begin{theorem}[Upper bound]
$m(n) \le \lfloor \log_2 n \rfloor + \nu(n) - 1$, where $\nu(n)$ is the Hamming weight of $n$.
\end{theorem}

\begin{proof}
The binary method uses $\lfloor \log_2 n \rfloor$ doublings and $\nu(n) - 1$ additions (one per 1-bit excluding the leading bit).
\end{proof}

\begin{proposition}
The binary method is not always optimal: $m(15) = 5$ but the binary method uses $3 + 4 - 1 = 6$ steps.
\end{proposition}

\begin{proof}
The chain $1, 2, 3, 6, 12, 15$ has length 5. Exhaustive search confirms no chain of length 4 exists for $n = 15$.
\end{proof}

\begin{theorem}[IDDFS optimality]
Iterative deepening DFS with initial depth $\lceil \log_2 n \rceil$ finds a shortest addition chain.
\end{theorem}

\begin{proof}
IDDFS explores all chains of length $r$ before trying $r + 1$, and the search space is finite (chain elements $\le n$, strictly increasing). The first solution is therefore optimal.
\end{proof}

\begin{lemma}[Pruning criterion]
If the current chain has largest element $a_d$ and $r - d$ steps remain, then $a_d \cdot 2^{r-d} < n$ implies the branch cannot reach $n$.
\end{lemma}

\begin{proof}
Doubling $a_d$ at every remaining step yields at most $a_d \cdot 2^{r-d}$.
\end{proof}

\section{Algorithm}

\begin{enumerate}
\item For each $n = 2, \ldots, 200$, try depth limits $r = \lceil \log_2 n \rceil, \lceil \log_2 n \rceil + 1, \ldots$
\item At depth $d$ with chain $[a_0, \ldots, a_d]$, try $a_{d} + a_i$ for $i = d, d-1, \ldots, 0$.
\item Prune if $a_d \cdot 2^{r-d} < n$.
\item Record $m(n) = r$ when the first valid chain is found.
\item Return $\sum_{n=1}^{200} m(n)$.
\end{enumerate}

\section{Complexity Analysis}

\begin{itemize}
\item \textbf{Time:} Branching factor at depth $d$ is at most $d + 1$; maximum depth $r \le 11$. Pruning renders the search fast in practice.
\item \textbf{Space:} $O(r)$ for the DFS stack and chain.
\end{itemize}

\section{Answer}

\[
\boxed{1582}
\]

\end{document}
